\section{Explicit invariant calculations}\label{sec:scalar-certificates}

This appendix derives the orbit formula for the minimum-support holonomy
invariant and gives two explicit scalar LU separators.  The holonomy
calculation supports the pencil classification.  The scalar separators are
not used in the rigidity or logical-group results; they show what finite-copy
contractions can detect on exceptional strata despite the generic-constancy
result in Section~\ref{sec:invariants}.

\subsection{The holonomy orbit formula}
\label{sec:holonomy-certificate}

Let
\[
 \mathcal P=\{\{1,2\},\{3,4\},\{5,6\}\},\qquad
 W=\operatorname{Stab}_{S_6}(\mathcal P)\cong C_2\wr S_3.
\]
The three pairs in \(\mathcal P\) are visible in the displayed columns of
\(H_t\).  They reduce the apparent \(450\)-case calculation to ten orbit
representatives.

\begin{lemma}[Holonomy orbit formula]\label{lem:holonomy-orbit-formula}
On the regular pencil locus, the minimum-support holonomy multiset is
\begin{equation}\label{eq:holonomy-invariant-multiset}
 \{4^{[90]},\ (-1/z)^{[144]},\ ((2z+1)^2/z^2)^{[24]},\
   \lambda_+^{[96]},\lambda_-^{[96]}\},
\end{equation}
where multiplicity is written in brackets and \(\lambda_\pm\) are the two
roots, over an algebraic closure, of
\[
 X^2-8X+(8-16z-1/z).
\]
\end{lemma}

\begin{proof}
An unoriented holonomy loop is indexed by a pair of coordinates \(e\) and an
unordered pair of distinct four-supports containing \(e\).  Equivalently,
write it as \((ij;ab,cd)\), where \(ij=e\) and \(ab,cd\) are the two omitted
pairs, each disjoint from \(e\).  There are
\(15\binom62=225\) such loops.  The group \(W\) has ten orbits on them.
Subscripts below give their orbit sizes:
\[
\begin{aligned}
4:\quad&
 (12;56,34)_{3},\ (12;46,35)_{6},\
 (13;46,25)_{24},\ (13;24,56)_{12};\\
-1/z:\quad&
 (12;56,36)_{24},\ (13;46,45)_{24},\ (13;45,25)_{24};\\
(2z+1)^2/z^2:\quad& (12;46,45)_{12};\\
\lambda_+:\quad& (13;56,25)_{48};\\
\lambda_-:\quad& (13;24,46)_{48}.
\end{aligned}
\]
This orbit decomposition is purely combinatorial.  The incidences of
\(ij,ab,cd\) with the three blocks of \(\mathcal P\) distinguish the ten
representatives, and orbit--stabilizer gives the displayed sizes; their sum
is \(225\).

Substitution of the shortened \(X\)- and \(Z\)-generators into
\eqref{eq:minimum-support-holonomy} shows first, on generators of \(W\), that
the rational function \(r(M)\) is constant on each displayed orbit.
Evaluating one representative of each orbit then gives the displayed values.
For the final two representatives, their sum is
\(8\) and their product is \(8-16z-1/z\), which gives the stated quadratic.
Reversing the two four-supports replaces the holonomy by its inverse and
does not change \(r(M)=\Tr(M)^2/\det M\).  Each unoriented orbit is therefore
counted twice in the ordered \(15\cdot6\cdot5=450\) multiset.  The orbit sizes
give \(90,144,24,96,96\), proving \eqref{eq:holonomy-invariant-multiset}.
\end{proof}

The paper-local certificate independently replays the ten representative
substitutions and reconstructs the same multiset directly from the CSS
Lagrangian.

Because \(A,B\ne0\) on \(\mathcal U\), we have \(z\ne0\).  The two quadratic
roots in \eqref{eq:holonomy-invariant-multiset} are taken over an algebraic
closure; their symmetric multiset descends to \(\F_q\).  Before
specialization, each root has multiplicity \(96\).  Coincident values have
their multiplicities added.

The identity
\[
 4B(t)^2+A(t)^2=4G(t)^2
\]
shows that the excluded GRS locus is exactly \(z=-1/4\).  Away from it, the
constant \(4\) bin is isolated and the two quadratic roots are distinct.  The
bin containing \(-1/z\) can merge with neither, one, or both of the
weight-\(24\) contribution and one weight-\(96\) contribution.  Its resulting
multiplicity is \(144\), \(168\), \(240\), or \(264\).  No other bin has one
of these multiplicities.  Thus the multiset canonically recovers \(-1/z\),
including every modular collision.

\subsection{A marginal-moment separator}

For a four-party set \(T\), let \(\rho_T\) be the reduced state and embed
it into all six parties as
\[
 A_T=\rho_T\otimes I_{T^c}.
\]
For an unordered triple \(\{T,U,V\}\) of distinct four-sets, the number
\(\Tr(A_TA_UA_V)\) is invariant under arbitrary product unitaries; party
permutation merely permutes the 455 triples.  By
\eqref{eq:css-reduced-state},
\begin{equation}\label{eq:marginal-moment-rank}
 \Tr(A_TA_UA_V)=q^{-\rank(L_C(T)+L_C(U)+L_C(V))}.
\end{equation}
Here and below, ranks of label spaces and matching systems are over
\(\F_q\).
Indeed, the three stabilizer sums have total normalization \(q^{-12}\).
Only triples of labels with sum zero survive the full trace; isotropy of
\(L_C\) removes multiplication phases, and the kernel of the addition map
from the three two-planes has size \(q^{6-r}\), where \(r\) is the
\(\F_q\)-rank in \eqref{eq:marginal-moment-rank}.

Index a four-set by its omitted pair, an edge of \(K_6\).  The rank in
\eqref{eq:marginal-moment-rank} is four exactly when the three corresponding chord lines
concur in both the six-arc \(\mathcal A\) and its Gale dual
\(\mathcal A^*\), represented by any kernel basis of the arc matrix.
To see this, split each shortened CSS plane into its one-dimensional
\(X(C)\) and \(Z(C^\perp)\) parts.  For three omitted edges, the three
shortened \(X\)-lines span dimension two precisely when the corresponding
chords of \(\mathcal A^*\) concur; the \(Z\)-lines give the same criterion
for \(\mathcal A\).  Each span has dimension at least two, so their sum
has rank four exactly when both concurrences hold.

Sixty triples are stars and concur universally.  Any other three-edge
graph that contains two adjacent edges but is not a star has two chords
meeting at their common arc point and a third chord missing it, since the
columns form an arc.  The only remaining graph type is a perfect
matching.  Hence
\begin{equation}\label{eq:marginal-moment-count}
\#\{q^{-4}\text{ moments}\}=60+b(\mathcal A,\mathcal A^*),
\end{equation}
where \(b\) counts perfect matchings concurrent in both configurations.

\begin{lemma}[Concurrent matchings on a conic]\label{lem:conic-matchings}
Let \(S\) be six distinct points on a nonsingular conic over a finite
field \(k\) of odd characteristic other than \(5\).  At most six perfect
matchings of \(S\) have concurrent joining chords.
\end{lemma}

\begin{proof}
Identify the conic with \(\mathbb P^1\).  Projection from a point off the
conic pairs the points in every fiber and hence induces a projective
involution; conversely, the chords joining the three two-element orbits
of a projective involution meet at its center.  Thus the concurrent
matchings are exactly the fixed-point-free involutions in
\[
 H=\operatorname{Stab}_{\mathrm{PGL}_2(k)}(S).
\]
The correspondence is injective because an element of
\(\mathrm{PGL}_2(k)\) is determined by its action on three points.
The group \(H\) acts freely on ordered triples of distinct points of
\(S\), so \(\lvert H\rvert\mid6\cdot5\cdot4=120\).

Dickson's classification, including the \(p\)-regular cases
\cite[\S\S239--261]{Dickson1901}, and Faber's \(p\)-irregular
classification \cite[Theorem~1.1]{Faber2023} are used only to obtain a
uniform upper bound of six; no classification of the six-sets is needed.
The possible cyclic, dihedral, \(A_4\), \(S_4\), and faithful degree-six
\(A_5\) cases contribute at most \(1,6,3,6,0\) fixed-point-free
involutions, respectively.  For \(A_5\), each involution fixes two of the
six points.  In the dihedral case, if the rotation order exceeded six, every orbit away
from the two rotation-fixed points would have size at least that order,
so no invariant six-set could exist; for rotation order at most six the
two reflection classes together contribute at most six fixed-point-free
elements (and the possible central half-turn contributes only when one
reflection class has fixed points).  In
the characteristic-three semi-elementary case, the normal \(3\)-group
has order \(3\), while its cyclic complement embeds in
\(\operatorname{Aut}(C_3)\), giving only \(C_3\) or \(S_3\).
The subfield groups whose order divides \(120\) are
\(\mathrm{PSL}_2(3)\cong A_4\) and
\(\mathrm{PGL}_2(3)\cong S_4\), already listed.  The remaining
characteristic-five subfield possibilities are excluded by the
hypothesis.  Thus every case is bounded by six.
\end{proof}

\begin{theorem}[Uniform \(H_3\)/GRS separation]\label{thm:lu-h3-grs}
Let \(\tau^2=\tau+1\), and let \(\mathfrak p\) be a prime ideal of
\(\mathbb Z[\tau]\) of odd residue characteristic for which the reduction
of the integral \(H_3\) (icosahedral) sextuple
\[
((0,1,1-\tau),(0,1,\tau-1),(1,1-\tau,0),(1,\tau-1,0),
 (1,0,-\tau),(1,0,\tau))
\]
is a six-arc.  Write \(k=\mathbb Z[\tau]/\mathfrak p\).  If the reduced
arc is non-GRS, put \(q=\lvert k\rvert\).  Its
equal-phase state is not LU-equivalent, even after a tensor-factor permutation, to
any six-point GRS state over \(k\).  (The hypothesis is vacuous in
characteristic five, where the reduction is itself GRS.)
\end{theorem}

\begin{proof}
Exact integral chord determinants give ten common concurrent matchings
for the \(H_3\) arc and its Gale dual.  The other five determinant norms are
\(-64\) and \(-4\), so the count is exactly ten in every residue
reduction covered by the theorem.  This is an exact fifteen-matching
calculation in \(\mathbb Z[\tau]\), not a field sample.  The \(H_3\) moment
multiset therefore contains \(70\) copies of \(q^{-4}\).

Lemma~\ref{lem:conic-matchings} gives at most six concurrent matchings
on a conic.  The bound is attained whenever \(\mathrm{PGL}_2(q)\) contains
an octahedral \(S_4\), by the corresponding six-set: the six
half-turns about axes through pairs of opposite edges act without fixed
vertices.  Thus a GRS state has at most \(66\) copies of \(q^{-4}\), and
the margin \(70>66\) is sharp for this argument.  In
characteristic five no subgroup bound is needed:
\(X^2-X-1=(X-3)^2\), so the residue of \(\tau\) is \(3\), and direct
substitution gives \(G(\tau)=0\).  The \(H_3\) reduction is therefore GRS,
contrary to the theorem's non-GRS hypothesis.
Equation~\eqref{eq:marginal-moment-rank} is
LU-invariant, so the gap proves the claim.
\end{proof}

The five nonconcurrent \(H_3\) matchings form a one-factorization of \(K_6\),
that is, a partition of its fifteen edges into five perfect matchings.
The marginal incidence recovers its \(S_5\) stabilizer but forgets the
even \(A_5\) orientation.  This is necessary: full local Fourier
transform sends \(\Psi_C\) to \(\Psi_{C^\perp}\), and an integral
diagonal equivalence to the dual returns the \(H_3\) code through an odd pentad permutation in every
odd reduction.  The orientation records which CSS summand is designated
\(X\) and which is designated \(Z\); local Clifford equivalence forgets this
choice, so it does not define a separate local-unitary class.

\subsection{The exact dimension-thirteen four-copy witness}

Use the permutation contraction
\(I_{\boldsymbol\sigma}\) and matching system
\(M_{\boldsymbol\sigma}(G)\) from
\eqref{eq:permutation-contraction}.

\begin{theorem}[Exact four-copy separator at \(q=13\)]
\label{thm:q13-lu}
The two regular \(q=13\) pencil classes with \(z=4\) and \(z=12\) are
not local-unitary equivalent, even after an arbitrary tensor-factor permutation.
\end{theorem}

\begin{proof}
At \(q=13\), the regular non-GRS pencil has two classes in the same complete
three-copy and marginal-moment bucket, with \(z=4\) and \(z=12\).
Use the explicit four-copy pattern
\[
\boldsymbol\sigma=(1,(03)(12),(23),(02)(13),(01)(23),(021)).
\]
For \(\pi\in S_6\), put
\((\pi\boldsymbol\sigma)_i=\sigma_{\pi^{-1}(i)}\), and sum
\eqref{eq:permutation-contraction}
over all tensor-factor permutations:
\[
 J_{\boldsymbol\sigma}(\Psi)
   =\sum_{\pi\in S_6}I_{\pi\boldsymbol\sigma}(\Psi).
\]
In \(\F_{13}\), one has \(4/9=12\), while \(4\) lies on neither component
of the transport divisor.  Theorem~\ref{thm:transport-divisor} therefore
gives rank \(21\) for all \(720\) terms at \(z=4\), and ranks \(20\) and
\(21\) with multiplicities \(192\) and \(528\) at \(z=12\).
Therefore
\begin{equation}\label{eq:q13-values}
 J_{\boldsymbol\sigma}(z=4)=720\cdot13^{-9},\qquad
 J_{\boldsymbol\sigma}(z=12)=3024\cdot13^{-9}.
\end{equation}
This proves arbitrary-LU inequivalence directly.  The complete two- and
three-copy contraction sets do not separate this pair;
\eqref{eq:q13-values} makes no
claim that four copies are globally minimal.
\end{proof}

Over \(\mathbb Q(t)\), every fixed four-copy matching matrix has generically
constant rank, so \eqref{eq:q13-values} detects a rank-jump stratum rather
than a generic orbit coordinate.  Section~\ref{sec:invariants} gives the
general distinction between these scalar contractions and the
operator-valued reduced states used in the rigidity theorem.
